\documentclass[10pt]{article}
\usepackage[T1]{fontenc}
\usepackage[utf8]{inputenc}
\usepackage[french]{babel}
\usepackage{amsmath,amssymb,mathrsfs,mathpazo}
\usepackage{geometry}
\geometry{a4paper,margin=2cm}
\usepackage{graphicx,array,tabularx,tabularray,colortbl,multirow,multicol,booktabs,enumitem,pifont,ulem,fancybox,diagbox,cancel}
\usepackage{pstricks,pstricks-add,pst-plot,pst-tree,pst-node,pst-text,pst-func,pst-math,pst-bspline,pst-3dplot,pst-eucl,pst-solides3d}
\let\pscalculate\undefined
\usepackage{pst-bezier}
\usepackage{tikz,pgfplots,tkz-tab}
\usetikzlibrary{arrows,arrows.meta,patterns,calc,positioning,trees,decorations.markings,shapes.geometric,shapes.misc,shapes.arrows,intersections,decorations.pathmorphing,decorations.pathreplacing,backgrounds,shadows,babel}
\pgfplotsset{compat=1.18}
\usepgfplotslibrary{fillbetween}
\usepackage[np]{numprint}
\usepackage{listings,makecell,pas-tableur,twoopt}
\let\barre\undefined
\usepackage[pstricks]{tablvar}
\definecolor{darkgreen}{RGB}{0,100,0}
\providecommand{\red}{\color{red}}
\providecommand{\blue}{\color{blue}}
\providecommand{\green}{\color{green}}
\providecommand{\black}{\color{black}}
\providecommand{\white}{\color{white}}
\providecommand{\gray}{\color{gray}}
\providecommand{\cyan}{\color{cyan}}
\providecommand{\magenta}{\color{magenta}}
\providecommand{\yellow}{\color{yellow}}
\providecommand{\violet}{\color{violet}}
\providecommand{\orange}{\color{orange}}
\providecommand{\euro}{\texteuro}
\providecommand{\EUR}{\texteuro}
\providecommand{\ieme}{\textsuperscript{e}}
\providecommand{\degres}{\ensuremath{^\circ}}
\providecommand{\degree}{\ensuremath{^\circ}}
\providecommand{\vv}[1]{\overrightarrow{#1}}
\providecommand{\Coord}[1]{#1}
\providecommand{\decosix}{}
\providecommand{\starredbullet}{\ensuremath{\bullet}}
\pagestyle{empty}
\setlength{\parindent}{0pt}
\usepackage{adjustbox}\DeclareUnicodeCharacter{2236}{:}\DeclareUnicodeCharacter{2190}{\ensuremath{\leftarrow}}
\providecommand{\R}{\mathbb{R}}
\providecommand{\N}{\mathbb{N}}
\providecommand{\Z}{\mathbb{Z}}
\providecommand{\Q}{\mathbb{Q}}
\providecommand{\D}{\mathbb{D}}
\providecommand{\C}{\mathbb{C}}
\providecommand{\e}{{\mathrm{e}}}
\providecommand{\vect}[1]{\overrightarrow{#1}}
\providecommand{\vectt}[1]{\overrightarrow{\mathrm{#1}}}
\providecommand{\vvt}[1]{\overrightarrow{\mathrm{#1}}}
\providecommand{\np}[2][]{{#2~#1}}
\providecommand{\vv}[1]{{\overrightarrow{#1}}}
\renewcommand{\d}{{\mathrm{d}}}
\providecommand{\degres}{{^\circ}}
\providecommand{\degre}{{^\circ}}
\providecommand{\blue}{\color{blue}}
\providecommand{\red}{\color{red}}
\providecommand{\enskip}{\quad}
\providecommand{\llbracket}{\lbrack\!\lbrack}
\providecommand{\rrbracket}{\rbrack\!\rbrack}
\providecommand{\euro}{{\text{€}}}
\providecommand{\no}{{\text{n°}}}
\providecommand{\interval}[2]{{[#1;#2]}}
\providecommand{\intervFO}[2]{{[#1;#2[}}
\providecommand{\dotfill}{\ldots}

\providecommand{\ds}{\displaystyle}
\providecommand{\textding}[1]{#1}
\providecommand{\Oij}{\(\mathrm{O};\vec{\imath},\vec{\jmath}\)}
\providecommand{\Oijk}{\(\mathrm{O};\vec{\imath},\vec{\jmath},\vec{k}\)}
\providecommand{\Ouv}{\(\mathrm{O};\vec{u},\vec{v}\)}
\providecommand{\intervFF}[2]{\ensuremath{[#1;#2]}}
\renewcommand{\euro}{\ifmmode\text{\texteuro}\else\texteuro\fi}
\DeclareUnicodeCharacter{00B0}{\ifmmode^{\circ}\else\textdegree\fi}
\DeclareUnicodeCharacter{20AC}{\ifmmode\text{\texteuro}\else\texteuro\fi}
% Énoncés : Éducation nationale. Transcription et corrigés : APMEP, Denis Vergès et contributeurs. Découpage et classement : CoopMaths, Sébastien Lozano, Éric Elter et contributeurs. Adaptation de la présentation : Maths Explorer.
% CC BY-SA : https://coopmaths.fr/www/about/
% https://coopmaths.fr/alea/static/bac/2023/tex/bac_2023_03_sujet1_metropole_3.tex
\begin{document}
\noindent\textit{APMEP / CoopMaths. CC BY-SA. Présentation adaptée par Maths Explorer.}\par\medskip
\medskip

Une entreprise a créé une Foire Aux Questions (\og FAQ \fg) sur son site internet.

\medskip

On étudie le nombre de questions qui y sont posées chaque mois.

\bigskip

\textbf{Partie A : Première modélisation}

\medskip

Dans cette partie, on admet que, chaque mois :

\begin{itemize}
\item[$\bullet~~$]90\,\% des questions déjà posées le mois précédent sont conservées sur la FAQ ;
\item[$\bullet~~$]130 nouvelles questions sont ajoutées à la FAQ.
\end{itemize}

Au cours du premier mois, $300$ questions ont été posées.

\medskip

Pour estimer le nombre de questions, en centaines, présentes sur la FAQ le $n$-ième
mois, on modélise la situation ci-dessus à l'aide de la suite $\left(u_n\right)$ définie par : \begin{center}$u_1 =3$ \quad et, pour tout entier naturel $n \geqslant 1,\:u_{n+1} = 0,9u_n + 1,3$.\end{center}

\smallskip

\begin{enumerate}
\item Calculer $u_2$ et $u_3$ et proposer une interprétation dans le contexte de l'exercice. \item Montrer par récurrence que pour tout entier naturel $n \geqslant 1$ :
\[u_n = 13 - \dfrac{100}{9} \times 0,9^n.\]
\end{enumerate}

\begin{minipage}{0.62\linewidth}
\begin{enumerate}[resume]
\item  En déduire que la suite $\left(u_n\right)$ est croissante.
\item On considère le programme ci-contre, écrit en langage Python.

Déterminer la valeur renvoyée par la saisie de seuil(8.5) et l'interpréter dans le contexte de l'exercice.
\end{enumerate}
\end{minipage}\hfill
\begin{minipage}{0.32\linewidth}
\renewcommand\arraystretch{0.9}
\begin{tabularx}{\linewidth}{|X|} \hline
def seuil(p) :\\
\qquad n=1\\
\qquad u=3\\
\qquad while u<=p :\\
\qquad  \qquad n=n+1\\
\qquad \qquad u=0.9*u+1.3 \\
\qquad return n\\ \hline
\end{tabularx}
\end{minipage}

\bigskip

\textbf{Partie B : Une autre modélisation}

\medskip

Dans cette partie, on considère une seconde modélisation à l'aide d'une nouvelle
suite $\left(v_n\right)$ définie pour tout entier naturel $n \geqslant 1$ par:
\[v_n = 9 - 6 \times \text{e}^{-0,19\times(n - 1)}.\]

Le terme $v_n$ est une estimation du nombre de questions, en centaines, présentes le $n$-ième mois sur la FAQ.

\medskip

\begin{enumerate}
\item Préciser les valeurs arrondies au centième de $v_1$ et $v_2$.
\item Déterminer, en justifiant la réponse, la plus petite valeur de $n$ telle que $v_n > 8,5$.
\end{enumerate}

\bigskip

\textbf{Partie C : Comparaison des deux modèles}

\medskip

\begin{enumerate}
\item L'entreprise considère qu'elle doit modifier la présentation de son site lorsque plus de 850 questions sont présentes sur la FAQ. 

Parmi ces deux modélisations, laquelle conduit à procéder le plus tôt à cette modification ?

Justifier votre réponse.
\item En justifiant la réponse, pour quelle modélisation y a-t-il le plus grand nombre de questions sur la FAQ à long terme?
\end{enumerate}

\bigskip

\end{document}
